\section{The boundary and the compact range}\label{bc:section}

We prove in this section the assertion of the main theorem whenever the
modulus of the proposed exceptional point is at least $1/10$.  Some of
the moment estimates obtained along the way will also be useful in the
remaining range.  Throughout the section, $P$ is monic of degree $n\ge6251$,
$N=n-1$, and all its zeros $z_1,\ldots,z_n$ belong to $\overline\D$.
Its critical points, with multiplicity, are denoted by
$\xi_1,\ldots,\xi_N$.  We put
\begin{equation}\label{bc:moments-definition}
 \mu=\frac1n\sum_i z_i=\frac1N\sum_j\xi_j,\qquad
 Q=\frac1N\sum_j|\xi_j|^2,\qquad
 L(z)=\frac1N\sum_j\log|z-\xi_j|.
\end{equation}
The letter $Q$ in this section denotes a second moment; reciprocal
energy will always be written $E_b$.

Suppose, towards a contradiction, that a real point $0<a\leq1$ in the
convex hull of the zeros satisfies
\begin{equation}\label{bc:hole}
 |a-\xi_j|>1\quad(1\leq j\leq N).
\end{equation}
Choose an actual zero $b=x+iy$ of largest real part, so $x\geq a$;
reflection permits us to assume $y\geq0$.  Write
\begin{equation}\label{bc:local-notation}
 r=|b-a|,\quad d=\min_j|b-\xi_j|,\quad
 S=\sum_j\frac1{b-\xi_j},\quad H=|S|,\quad
 E_b=\sum_j\frac1{|b-\xi_j|^2}.
\end{equation}
Gauss--Lucas and \eqref{bc:hole} imply
\begin{equation}\label{bc:entrance}
 |\xi_j|\leq1,\quad \Re\xi_j<a/2,\quad
 r^2\leq1-a^2<1,\quad d>1-r,\quad
 \Re S>\frac a2E_b.
\end{equation}
In particular $b$ is simple.  We use the actual-zero energy inequality
$E_b\geq N$ of Proposition~\ref{ar:energy}.  It excludes $r=0$ and also excludes a
strict witness which is itself an original zero.  Consequently all
applications below have $0<a<1$ and $r>0$.

\subsection{A product inequality retaining the original zeros}

The fiber arguments use the following fact for arbitrary equal-value
fibers.

\begin{lemma}[Equal-value perpendicular bisector]\label{bc:equal-fiber}
If $P(v)=P(w)$ and $v\ne w$, the critical points of $P$ cannot all
belong to either one of the two open half-planes bounded by the
perpendicular bisector of $[v,w]$.
\end{lemma}

\begin{proof}
We prove the needed half-plane polarization fact first.  Let $f$ be a
nonzero polynomial of degree at most $m$ with no zero in
$\{\Re z>0\}$.  For $\Re w>0$ put
\[
 T_w f(z)=f(z)+\frac{w-z}{m}f'(z).
\]
This polynomial has degree at most $m-1$ and is nonzero throughout
$\{\Re z>0\}$.  To see the latter assertion, write the zeros of
$f$ as $\zeta_1,\ldots,\zeta_d$, where $d\leq m$, and fix
$z$ with $x=\Re z>0$.  Each $v_j=(z-\zeta_j)^{-1}$ satisfies
$\Re v_j\geq x|v_j|^2$.  The disk
$\{v:\Re v\geq x|v|^2\}$ is convex and contains zero; hence
$v=m^{-1}\sum_{j=1}^d v_j$ is in that disk as well.  If $v=0$,
then $T_w f(z)=f(z)\ne0$.  Otherwise $T_w f(z)=0$ would imply
$w=z-1/v$, whose real part is at most zero, a contradiction.

Let $F(\zeta_1,\ldots,\zeta_m)$ be symmetric and affine in each
variable separately.  If $f(z)=F(z,\ldots,z)$ has no zero in the
right half-plane, neither does $F$ when all its variables are in that
half-plane.  Indeed, replacing one of the remaining equal variables
$z$ by $w$ applies exactly the operation $T_w$, with $m$ replaced by
the number of remaining variables.  Iteration proves the assertion.
Reflection proves its left-half-plane version.

An affine change of variable reduces the asserted bisector to the
imaginary axis and the two equal-value points to $-1,1$.  Consider
\[
 F(\zeta_1,\ldots,\zeta_N)
   =\int_{-1}^{1}\prod_{j=1}^N(t-\zeta_j)\,dt.
\]
Its diagonal is
\[
 F(z,\ldots,z)=\frac{(1-z)^{N+1}-(-1-z)^{N+1}}{N+1}.
\]
Every zero of this polynomial has $|1-z|=|-1-z|$, and hence is on
the imaginary axis.  Half-plane polarization therefore makes $F$
nonzero if all the critical points are strictly on the same side of
that axis.  On the other hand, $F(\xi_1,\ldots,\xi_N)$ is a nonzero
constant multiple of $P(1)-P(-1)=0$.  This contradiction proves the
lemma.
\end{proof}

For $|\zeta|<1$ define
\[
 \mathcal B_P(\zeta)=
 \prod_{i=1}^n\left|\frac{\zeta-z_i}{1-\overline{z_i}\zeta}\right|.
\]
Every factor is at most one.  The following lower bound for the product
retains the reciprocal energy $E_b$ at the original zero $b$.

\begin{lemma}[Adaptive product inequality]\label{bc:blaschke}
Under \eqref{bc:hole}--\eqref{bc:local-notation}, put
\begin{equation}\label{bc:adaptive-def}
 u=\min(r^2,1/2),\qquad t=1-u,\qquad
 \gamma=b-ta,\qquad C=ta\bigl(2x-(1+t)a\bigr).
\end{equation}
Then $|\gamma|<1$, $C\geq a^2u(1-u)>0$, and
\begin{equation}\label{bc:blaschke-bound}
 C\exp(CE_b/2)
 <2(1+|b|)(1+|\gamma|)\mathcal B_P(\gamma)
 \leq8\mathcal B_P(\gamma).
\end{equation}
In particular, if $n\geq6251$, then
\begin{equation}\label{bc:Ccap}
 C<C_*:=\frac1{388}.
\end{equation}
\end{lemma}

\begin{proof}
Set $k=1-|\gamma|^2$.  Since
$|\gamma|^2=r^2+2ua(x-a)+u^2a^2>u$ and
$|b|^2-|\gamma|^2=ta(2x-ta)>0$, we have $0<k<t<1$.
Direct expansion gives, for every critical point,
\begin{align}
 |b-\xi|^2-|1-\overline\gamma\xi|^2
 &=t\bigl(|\xi-a|^2-1\bigr)
   +C+(t-k)(1-|\xi|^2)>C.             \label{bc:pointwise-product}
\end{align}
Also $|1-\overline\gamma\xi|^2-|\gamma-\xi|^2
=k(1-|\xi|^2)\geq0$.

Let $M=1+|b|$ and $\eta=C/(2M)$.  For every $z\in D(b,\eta)$,
\[
 |z-\xi|^2>(|b-\xi|-\eta)^2
 \geq |b-\xi|^2-2M\eta+\eta^2>|\gamma-\xi|^2.
\]
Here $\eta<|b-\xi|/2$, by \eqref{bc:pointwise-product}, and
$\eta<ta=|b-\gamma|$.  Lemma~\ref{bc:equal-fiber}
therefore forbids two equal-value points, one in $D(b,\eta)$ and the
other equal to $\gamma$.  In particular $P(\gamma)\ne0$.
If $|P(\gamma)|<\eta|P'(b)|/n$, the logarithmic derivative of
$P(z)-P(\gamma)$ at $b$ would give
\[
 \left|\sum_{P(\zeta)=P(\gamma)}\frac1{b-\zeta}\right|
 =\frac{|P'(b)|}{|P(\gamma)|}>\frac n\eta.
\]
At least one of these $n$ zeros, counted with multiplicity, would be in
$D(b,\eta)$, again a contradiction.  Hence
\begin{equation}\label{bc:fiber-lower}
 |P(\gamma)|\geq\frac{\eta}{n}|P'(b)|.
\end{equation}

Let $q=P'/n$, and introduce the reciprocal polynomials
\[
 P^*(z)=\prod_i(1-\overline{z_i}z),\qquad
 q^*(z)=\prod_j(1-\overline{\xi_j}z).
\]
From \eqref{bc:pointwise-product} and $-\log(1-v)>v$ for $0<v<1$,
\begin{equation}\label{bc:critical-product}
 \frac{|P'(b)|}{n|q^*(\gamma)|}>
 \exp\left(\frac C2\sum_j|b-\xi_j|^{-2}\right).
\end{equation}
The exact reversed derivative identity is
\[
 \frac{q^*(\gamma)}{P^*(\gamma)}
 =\frac1n\sum_i\frac1{1-\overline{z_i}\gamma}.
\]
For $|w|\leq g<1$ one has $\Re(1-w)^{-1}\geq(1+g)^{-1}$,
as follows by multiplying by $|1-w|^2$.  Thus the modulus of the last
display is at least $(1+|\gamma|)^{-1}$.  Combining it with
\eqref{bc:fiber-lower} and \eqref{bc:critical-product} proves
\eqref{bc:blaschke-bound}.  Finally
\[
 \frac{6250}{776}>\log(8\cdot388)
\]
implies \eqref{bc:Ccap}, since $C\exp(NC/2)$ increases strictly in
$C>0$ and in $N$.  This numerical comparison is certified by the
positive exponential series, as described below.
\end{proof}

The normalization at the origin uses the following movement estimate.

\begin{lemma}[Movement of an equal-value fiber]\label{bc:movement}
Let $\omega=b/|b|$, let $a_0\leq a$, and suppose
$2\Re(\overline\omega\xi_j)<d_0$ for all $j$.  If positive numbers
$s,\delta$ satisfy
\begin{equation}\label{bc:movement-gates}
 a_0-s>\delta,\qquad
 (a_0-s)(a_0+s-d_0)>\delta(4+\delta),
\end{equation}
then $|P(s\omega)|\geq\delta|P'(b)|/n$.
\end{lemma}

\begin{proof}
Writing $v=|b|\geq a_0$, the squared-distance difference at $b$ is
\[
 |b-\xi|^2-|s\omega-\xi|^2
 =(v-s)\bigl(v+s-2\Re(\overline\omega\xi)\bigr)
 >(a_0-s)(a_0+s-d_0).
\]
Moving $b$ by at most $\delta$ decreases this difference by at most
$\delta(4+\delta)$.  Thus every point of $D(b,\delta)$ is farther
from every critical point than $s\omega$ is.  The first condition in
\eqref{bc:movement-gates} puts $s\omega$ outside that disk.  The same
fiber and logarithmic-derivative argument as in
\eqref{bc:fiber-lower} proves the assertion.
\end{proof}

\subsection{Normalization, annuli, and moment estimates}

\begin{proposition}[Three moment estimates]\label{bc:moment-proposition}
Assume $n\geq6251$.  Under \eqref{bc:hole}, each row of
Table~\ref{bc:moment-conclusions} implies $P(0)\ne0$.  With
\begin{equation}\label{bc:B-and-error}
 B=-\frac1n\log|P(0)|,\qquad
 \mathcal E(r;a,N)=\frac{8r}{aN}
       +\frac{r^2}{2(1-r)(1-2r)},
\end{equation}
its conclusions are
\begin{align}
 B&<\frac\ell n+\mathcal E(r;a,N),                 \label{bc:moment-B}\\
 |\mu|&\leq\kappa B+n^{-2},                       \label{bc:moment-mean}\\
 Q&<1-a^2+w\left(\frac\ell N+\mathcal E(r;a,N)\right)
                   +2n^{-2}.                    \label{bc:moment-Q}
\end{align}
\begin{table}[ht]
\centering
\caption{The moment estimates.  The hypotheses use the actual $a,r$.}
\label{bc:moment-conclusions}
\begin{tabular}{@{}cccccc@{}}\toprule
Row&$a\geq$&$r\leq$&$\ell$&$\kappa$&$w$\\\midrule
M0&$.9$&$.057$&$4.42$&$2.5482$&$5.098$\\
M1&$.93$&$.055$&$4.22$&$2.3817$&$4.765$\\
M2&$.95$&$.054$&$4.11$&$2.3108$&$4.623$\\\bottomrule
\end{tabular}
\end{table}
The same formulas hold with the parameters in the four boundary rows
of the tables below, whenever their respective hypotheses hold.
\end{proposition}

We give the normalization argument with parameters; this also supplies
the estimates needed to prove the boundary assertions.  A row specifies
lower $a_0$, upper displacement $e<1/2$, a squared critical-distance
bound $R_c^2<(1-e)^2$, and the constants listed in
Tables~\ref{bc:preliminary-table}--\ref{bc:annulus-table}.
The names B, A, G1, G2 stand, respectively, for the boundary range
$a\geq.975$, and the small displacement ranges
$(a,r)\in[.85,1)\times[0,.02]$,
$[.96,1)\times[0,.04]$, and
$[.965,1)\times[0,.045]$.
For row B the displacement bound is a conclusion of
Lemma~\ref{bc:blaschke}: indeed
\[
 (.975)^2(.053)^2\bigl(1-(.053)^2\bigr)>1/388,
\]
and $u(1-u)$ increases on $[0,1/2]$; hence $r<.053$.

\begin{proof}[Proof of Proposition~\ref{bc:moment-proposition}]
First,
\[
 |\omega-1|^2=\frac{2y^2}{|b|(|b|+x)}\leq r^2/a^2,
 \qquad
 2\Re(\overline\omega\xi_j)<a+2r/a.
\]
All rows satisfy $2e<a_0^2$.  The function $a+2e/a$ therefore
increases on $[a_0,1]$, and is at most $1+2e<d_0$ there.
The two pairs $(s_1,\delta_1),(s_2,\delta_2)$ in the tables satisfy
\eqref{bc:movement-gates}.  Lemma~\ref{bc:movement} applies to both.

Put $k=N/n$, $A_n=\delta_1^{2/n}$, and
$A_-=1-\alpha_0/n$.  The inequality
$\alpha_0/2>\log(1/\delta_1)$ gives $A_n>A_->0$.
Original-zero AM--GM at $s_1\omega$ gives
\begin{equation}\label{bc:directional-mean}
 2s_1\Re(\overline\omega\mu)
 \leq1+s_1^2-A_n e^{2kL(b)}.
\end{equation}
On $R_c\leq s\leq2$ the logarithmic chord
\begin{equation}\label{bc:coarse-chord}
 s^2\leq R_c^2+\frac{22}5\log(s/R_c)
\end{equation}
follows by convexity of $e^{2v}$ in $v=\log s$ and its endpoint
inequality.  The latter follows already from
$\log(1/R_c)\geq(1-R_c^2)/2$ and the displayed certificate below.
Average \eqref{bc:coarse-chord} at $s=|b-\xi_j|$ and use
\eqref{bc:directional-mean}.  For $v=|b|$ this yields
\begin{equation}\label{bc:F-definition}
 Q\leq F(L(b)):=
 \frac v{s_1}(1+s_1^2-A_ne^{2kL(b)})+R_c^2-v^2
     +\frac{22}5(L(b)-\log R_c).
\end{equation}
Let
\[
 T=\frac{1+s_1^2-A_-R_c^2}{s_1},\qquad
 F_0=a_0T+R_c^2-a_0^2.
\]
The row comparisons $T<2a_0$ and
$22/5-2a_0kA_-R_c^2/s_1<0$ show that $F$ decreases on the
entire half-line $L(b)\geq\log R_c$.  Indeed $F''<0$, while
$R_c^{2k}\geq R_c^2$; the quadratic $vT-v^2$ decreases for
$v\geq a_0$.  The resulting chain is
\begin{equation}\label{bc:coarse-correct-chain}
 Q\leq F(L(b))\leq F(\log R_c)\leq F_0<q_c,
 \qquad |\mu|<m_c,
\end{equation}
where the last inequality uses $q_c<m_c^2$.

The same calculation at every $z\omega$, $0\leq z\leq s_2$, bounds
the average squared critical distance by
\[
 z^2+\frac{v-z}{s_1}(1+s_1^2-A_ne^{2kL(b)})+R_c^2-v^2
       +\frac{22}5(L(b)-\log R_c).
\]
Subtract $.9e^{2L(b)}$.  Its derivative in $L(b)$ is negative
throughout the same half-line, because
\[
 2R_c^2\left(kA_-\frac{a_0-s_2}{s_1}+.9\right)>22/5.
\]
At $L(b)=\log R_c$ it is at most
$z^2-zT+F_0-.9R_c^2$.  This is a convex quadratic in $z$;
the two endpoint values, at $0,s_2$, are negative.
Critical-zero AM--GM consequently gives
\begin{equation}\label{bc:primitive-decay}
 |P'(z\omega)|<(.95)^N|P'(b)|\qquad(0\leq z\leq s_2).
\end{equation}
Since $e^{N/20}>200n$ and $s_2<100\delta_2$, integration and the
second movement inequality give
\begin{equation}\label{bc:normalization}
 |P(0)|\geq\frac{\delta_2}{2n}|P'(b)|>0,
 \qquad
 L(b)\leq\frac\ell N-\frac nN B,
\end{equation}
where $e^\ell>2/\delta_2$.  The inequalities
$R_c>1.001p$ and $(1.001)^N>2p/\delta_2$ also give
$|P(0)|>p^n$.  In particular
$0\leq B< -\log R_c+\ell/n<B_c$.

For the annulus iteration, suppose $|\mu|\leq m$,
$Q<q_c$, and
\begin{equation}\label{bc:annulus-gates}
 \rho^2+2\rho m+q_c<s_*^2,\qquad
 \rho<s_*<p,\qquad n(s_*/p)^n<1.
\end{equation}
For every $|z|\leq\rho$, AM--GM gives
$|P'(z)|<ns_*^{N}$.  Integrating on the segment from $0$ to $z$,
\[
 \frac{|P(z)-P(0)|}{|P(0)|}
 <\frac\rho{s_*}n(s_*/p)^n<1.
\]
Thus every original zero has modulus greater than $\rho$.
For $0<\rho\leq s\leq1$, the function
$2\sinh t/t$, $t=-\log s$, increases with $t\geq0$.
It follows that
\[
 s^{-1}-s\leq
 \frac{\rho^{-1}-\rho}{-\log\rho}(-\log s).
\]
Since $-P'(0)/P(0)=\sum_i z_i^{-1}$, any
$\kappa>(\rho^{-1}-\rho)/(-\log\rho)$ therefore gives
\begin{equation}\label{bc:reciprocal-mean}
 |\mu|\leq\kappa B+
 \frac{|P'(0)|}{n|P(0)|}.
\end{equation}
The residual in this formula is bounded by
\[
 \frac{m_c^N}{p^n}=\frac{(m_c/p)^n}{m_c}<2d_s^n<n^{-2},
\]
using $m_c/p<d_s$ and $1/m_c<2$.
Beginning with $m=m_c$,
each successive triple $(\rho,\kappa,m')$ in
Table~\ref{bc:annulus-table} satisfies \eqref{bc:annulus-gates}
and $\kappa B_c+n^{-2}<m'$.  The last triple supplies
\eqref{bc:moment-mean}.

The geometric bounds \eqref{bc:entrance} and $E_b\geq N$ give
\[
 H>Na/2,\qquad E_b/H<2/a,\qquad
 E_b/H^2<4/(Na_0^2),\qquad (Hd)^{-1}<2/[Na_0(1-e)].
\]
The last two bounds are respectively below $7/625$ and $1/270$.
The Fine companion estimate of Proposition~\ref{lt:fine200} therefore gives
\begin{equation}\label{bc:fine-phase}
 |\Im S|<4E_b/H<8/a.
\end{equation}
For $w_0=b-a$, the complete logarithmic series at $b$ converges
absolutely, because $r/d<r/(1-r)<1$.  Its linear term satisfies
\[
 -\Re(w_0S)=-(x-a)\Re S+y\Im S\leq r|\Im S|.
\]
The terms of order at least two have total absolute value, after
averaging, at most
$E_br^2/[2N(1-r/d)]$.  Using $E_b/N<(1-r)^{-2}$ proves
\begin{equation}\label{bc:reverse-log}
 L(a)\leq L(b)+\mathcal E(r;a,N).
\end{equation}
Since $L(a)>0$, \eqref{bc:normalization} proves
\eqref{bc:moment-B}.

Finally $s^2\leq1+w\log s$ on $1\leq s\leq2$, because
$w\log2>3$ and the exponential is convex.  Apply this chord at
$s=|a-\xi_j|$ and use
\eqref{bc:normalization}, \eqref{bc:reciprocal-mean}, and
\eqref{bc:reverse-log}.  We obtain
\[
 Q<1-a^2+w\left(\frac\ell N+\mathcal E(r;a,N)\right)
      +2n^{-2}+(2a\kappa-wn/N)B.
\]
The coefficient of $B$ is negative, since $w>2\kappa$ and $a\leq1$.
This proves \eqref{bc:moment-Q}.
\end{proof}


\begin{table}[ht]
\centering\small
\caption{Preliminary parameters.  In every row $s_1=.25$ and $s_2=.45$.}
\label{bc:preliminary-table}
\begin{tabular}{@{}crrrrrrrr@{}}\toprule
Row&$a_0$&$e$&$R_c^2$&$d_0$&$\delta_1$&$\delta_2$&$\alpha_0$&$\ell$\\\midrule
B & 0.975 & 0.053 & 0.896708 & 1.10611 & 0.01976 & 0.03838 & 7.86 & 3.97\\
A & 0.85 & 0.02 & 0.960299 & 1.04011 & 0.00823 & 0.02383 & 9.61 & 4.44\\
G1 & 0.96 & 0.04 & 0.921499 & 1.08011 & 0.02114 & 0.03857 & 7.73 & 3.96\\
G2 & 0.965 & 0.045 & 0.911924 & 1.09011 & 0.02047 & 0.03836 & 7.79 & 3.97\\
M0 & 0.9 & 0.057 & 0.889148 & 1.11411 & 0.00534 & 0.02433 & 10.48 & 4.42\\
M1 & 0.93 & 0.055 & 0.892924 & 1.11011 & 0.01089 & 0.0297 & 9.05 & 4.22\\
M2 & 0.95 & 0.054 & 0.894815 & 1.10811 & 0.01474 & 0.03346 & 8.45 & 4.11\\
\bottomrule\end{tabular}
\end{table}
\begin{table}[ht]
\centering\small
\caption{Coarse normalization and annulus parameters.}
\label{bc:coarse-table}
\begin{tabular}{@{}crrrrrrr@{}}\toprule
Row&$q_c$&$m_c$&$p$&$s_*$&$B_c$&$d_s$&$w$\\\midrule
B & 0.59717 & 0.77287 & 0.94595 & 0.94414 & 0.055158 & 0.81705 & 4.494\\
A & 0.59041 & 0.76849 & 0.97891 & 0.9771 & 0.020976 & 0.78506 & 4.33\\
G1 & 0.54582 & 0.7389 & 0.95893 & 0.95712 & 0.041521 & 0.77056 & 4.33\\
G2 & 0.56641 & 0.75271 & 0.95394 & 0.95213 & 0.046745 & 0.78907 & 4.354\\
M0 & 0.70869 & 0.84194 & 0.94195 & 0.94014 & 0.059463 & 0.89384 & 5.098\\
M1 & 0.66376 & 0.81482 & 0.94395 & 0.94214 & 0.057312 & 0.86322 & 4.765\\
M2 & 0.63422 & 0.79648 & 0.94495 & 0.94314 & 0.056237 & 0.8429 & 4.623\\
\bottomrule\end{tabular}
\end{table}
\begin{longtable}{@{}clll@{}}
\caption{Successive annulus triples $(\rho,\kappa,m')$.
A dash means that only the resulting inequality $|\mu|\leq\kappa B+n^{-2}$ is retained.}
\label{bc:annulus-table}\\
\toprule Row&$\rho$&$\kappa$&$m'$\\\midrule\endfirsthead
\toprule Row&$\rho$&$\kappa$&$m'$\\\midrule\endhead
B & 0.169 & 3.2333 & 0.178344\\
 & 0.39 & 2.3091 & 0.127367\\
 & 0.427 & 2.2504 & 0.124129\\
 & 0.43 & 2.2462 & 0.123897\\
 & 0.43 & 2.2462 & --\\
A & 0.206 & 2.9424 & 0.061721\\
 & 0.543 & 2.1268 & 0.044613\\
 & 0.558 & 2.1155 & 0.044376\\
 & 0.558 & 2.1155 & --\\
G1 & 0.216 & 2.8802 & 0.11959\\
 & 0.498 & 2.1661 & 0.08994\\
 & 0.523 & 2.1432 & 0.088989\\
 & 0.523 & 2.1432 & --\\
G2 & 0.197 & 3.0035 & 0.1404\\
 & 0.457 & 2.2109 & 0.10335\\
 & 0.486 & 2.1783 & 0.101826\\
 & 0.488 & 2.1762 & 0.101728\\
 & 0.488 & 2.1762 & --\\
M0 & 0.096 & 4.4043 & 0.261894\\
 & 0.229 & 2.8073 & 0.166932\\
 & 0.281 & 2.5822 & 0.153547\\
 & 0.29 & 2.5515 & 0.151721\\
 & 0.291 & 2.5482 & --\\
M1 & 0.125 & 3.7872 & 0.217054\\
 & 0.301 & 2.5165 & 0.144227\\
 & 0.348 & 2.3928 & 0.137138\\
 & 0.353 & 2.3817 & 0.136502\\
 & 0.353 & 2.3817 & --\\
M2 & 0.144 & 3.5093 & 0.197354\\
 & 0.343 & 2.4043 & 0.135212\\
 & 0.385 & 2.318 & 0.130359\\
 & 0.389 & 2.3108 & 0.129954\\
 & 0.389 & 2.3108 & --\\
\bottomrule\end{longtable}
\begin{table}[ht]
\centering\small
\caption{Uniform bounds $Q<q_0$, $B<B_0$, $|\mu|<m_0$, and
$|P(0)|>p_f^n$ furnished by the moment formulas at $n=6251$.}
\label{bc:refined-table}
\begin{tabular}{@{}crrrr@{}}\toprule
Row&$q_0$&$B_0$&$m_0$&$p_f$\\\midrule
B & 0.06001 & 0.0023642 & 0.005311 & 0.997625\\
A & 0.28164 & 0.0009535 & 0.0020177 & 0.999036\\
G1 & 0.08531 & 0.0015932 & 0.0034151 & 0.998396\\
G2 & 0.07689 & 0.0018604 & 0.0040492 & 0.998129\\
M0 & 0.20395 & 0.0027331 & 0.0069651 & 0.997256\\
M1 & 0.14726 & 0.0025497 & 0.0060732 & 0.99744\\
M2 & 0.10888 & 0.0024586 & 0.0056819 & 0.997531\\
\bottomrule\end{tabular}
\end{table}

The numerical statements used in the preceding proof are finite rational
comparisons, with logarithms bounded by convergent positive series.
The remaining comparisons are as follows, with $n_0=6251$,
$N_0=6250$, $k_0=N_0/n_0$, and $A_-=1-\alpha_0/n_0$:
\begin{align}
 &2e<a_0^2,\quad 1+2e<d_0,\quad R_c^2<(1-e)^2,
       \quad\eqref{bc:movement-gates}\text{ for }j=1,2;
                                                        \label{bc:pre-gates-1}\\
 &R_c^2+\tfrac{22}5\bigl(\log2+(1-R_c^2)/2\bigr)>4,
       \quad \alpha_0/2>\log(1/\delta_1),
       \quad e^\ell>2/\delta_2;                         \label{bc:pre-gates-2}\\
 &T<2a_0,\quad F_0<q_c<m_c^2,\quad
       2a_0k_0A_-R_c^2/s_1>22/5;                       \label{bc:pre-gates-3}\\
 &2R_c^2[k_0A_-(a_0-s_2)/s_1+.9]>22/5,
       \quad z^2-zT+F_0<.9R_c^2\quad(z=0,s_2);         \label{bc:pre-gates-4}\\
 &N_0/20>\log(200n_0),\quad s_2<100\delta_2,
       \quad R_c>1.001p,
       \quad N_0\log(1.001)>\log(2p/\delta_2);         \label{bc:pre-gates-5}\\
 &-\log R_c+\ell/n_0<B_c,\quad 1/m_c<2,
       \quad m_c/p<d_s,
       \quad (1-d_s)n_0>\log(2n_0^2);                 \label{bc:pre-gates-6}\\
 &-n_0\log(s_*/p)>\log n_0,\quad
       -n_0\log(s_*/p)>1,\quad(1-d_s)n_0>2;           \label{bc:pre-gates-7}\\
 &\rho^2+2\rho m+q_c<s_*^2,\quad \rho<s_*,\quad
       \rho^{-1}-\rho<-\kappa\log\rho,
       \quad \kappa B_c+n_0^{-2}<m';                 \label{bc:pre-gates-8}\\
 &\frac4{N_0a_0^2}<\frac7{625},\quad
       \frac2{N_0a_0(1-e)}<\frac1{270},\quad
       w>2\kappa,\quad w\log2>3.                     \label{bc:pre-gates-9}
\end{align}
In \eqref{bc:pre-gates-8}, $m=m_c$ at the first stage and is the
preceding $m'$ thereafter; the last inequality is omitted for a dash.
The caps in Table~\ref{bc:refined-table} are checked by substituting
$a=a_0,r=e,n=n_0$ in \eqref{bc:moment-B}--\eqref{bc:moment-Q}, and by
$\kappa B_0+n_0^{-2}<m_0$, $1-B_0>p_f$.
These conditions prove all the displayed estimates for every $n\geq n_0$:
$k,A_-$ increase, the moment errors decrease, and
$n(s_*/p)^n$ and $2n^2d_s^n$ decrease after their stated derivative
bounds.

\subsection{Support and the normal of a boundary edge}

The following transfer is a finite identity between the original and
critical empirical measures.  It retains the normalization difference
between $n$ and $N$.

\begin{lemma}[Moment support bound]\label{bc:poisson-support}
Suppose $P(0)\ne0$, $B\leq B_0$, $Q\leq q<1$, and $|\mu|\leq m$.
Let $0<h<1$ and put
\[
 \Lambda=\frac{(1-q)^2}{1+q+2m}.
\]
If
\begin{equation}\label{bc:support-test}
 (N\Lambda+1)(1-h)^2>2(1-h)+2nB_0,
\end{equation}
then the convex hull of the original zeros contains $\overline{D(0,h)}$
in its interior.
\end{lemma}

\begin{proof}
Rotate a proposed contrary support direction to the real axis, so that
$\Re z_i\leq h$ for all $i$.  Define
\[
 C_o(z)=\prod_i(1-z_i z),\quad C_c(z)=\prod_j(1-\xi_jz),
 \quad f(z)=\frac1n\sum_i\frac1{1-z_i z},\quad
 g(z)=\frac1N\sum_j\frac1{1-\xi_jz}.
\]
The exact identities $C_c=C_o-zC_o'/n$ and $C_c/C_o=f$ give
\begin{equation}\label{bc:poisson-identity}
 N(g-f)=f-1-zf'/f.
\end{equation}
The support assumption, and Gauss--Lucas for the critical points,
make every term well defined at $z=1$.  Also
\[
 |f'(1)|\leq\frac1n\sum_i|1-z_i|^{-2},\qquad
 \Re f(1)\geq(1-h)\frac1n\sum_i|1-z_i|^{-2}.
\]
Consequently $f(1)\ne0$ and $|f'(1)/f(1)|\leq(1-h)^{-1}$.
Set $p_o=2\Re f(1)-1$, $p_c=2\Re g(1)-1$.  Then
\[
 p_o=\frac1n\sum_i\frac{1-|z_i|^2}{|1-z_i|^2}
 \leq\frac{2B_0}{(1-h)^2},\qquad
 Np_c=np_o-1-2\Re(f'(1)/f(1)).
\]
Here we used $1-|z_i|^2\leq-2\log|z_i|$.
For the other measure, Cauchy--Schwarz gives
\[
 p_c\left(\frac1N\sum_j|1-\xi_j|^2\right)
 \geq\left(\frac1N\sum_j\sqrt{1-|\xi_j|^2}\right)^2
 \geq(1-Q)^2.
\]
Its denominator equals $1+Q-2\Re\mu$, whence $p_c\geq\Lambda$.
Combining these estimates contradicts \eqref{bc:support-test}.
Since every direction was allowed, all support values exceed $h$.
\end{proof}

We use the inverse-log quantities of Lemma~\ref{lt:inverse-log}.
Suppose, at the actual rightmost zero $b$,
\begin{equation}\label{bc:field-hypotheses}
 H>H_0>0,\quad E_b/H<V,\quad d^{-1}<M.
\end{equation}
For $R>0$ define
\begin{align}
 \lambda&=\frac{VR}{1-MR},&
 e_R&=V\left(\frac{R^2}2+\frac{MR^3}{3(1-MR)}\right),
 &\widehat r&=R-e_R,                                  \label{bc:inverse-def-1}\\
 B_2&=\frac{V}{(1-MR)^2(1-\lambda)^3},&
 p_-&=(1+\lambda)^{-1},&p_+&=(1-\lambda)^{-1}.
                                                               \label{bc:inverse-def-2}
\end{align}
Choose $0<\nu<1$, $W>0$, and $e_W>e^{-W}$, and put
\[
 \epsilon=1.001B_2/H_0,\qquad g=1.001\nu B_2\widehat r.
\]
The sufficient comparisons are
\begin{align}
 &MR<1,\quad 0<\lambda<1,\quad \widehat r>0,
       \quad W^2<(1-\nu^2)(H_0\widehat r)^2,
                                                    \label{bc:inverse-gates-1}\\
 &\epsilon(1+e_W)+ge_W<1-p_+e_W,
                                                    \label{bc:inverse-gates-2}\\
 &\epsilon(1+e_W)+ge_W^2<p_--e_W,
       \qquad \epsilon+ge_W<p_-.                   \label{bc:inverse-gates-3}
\end{align}
They ensure that the analytic logarithm
$\mathcal L(z)=\log(P'(z)/P'(b))$, fixed at $b$, is injective on
$D(b,R)$ and has an inverse $\phi$ on $D(0,H\widehat r)$.
Moreover
\begin{equation}\label{bc:inverse-location}
 |\phi(w)-b-w/S|<e_R,\qquad
 \frac{SP(\phi(w))}{P'(b)}=\int_0^we^v A(v)\,dv,\quad A=S\phi',
\end{equation}
with $\Re A>p_-$, $|A|<p_+$ and $|A'|<B_2/H$.
On every rectangle with vertical sides $\Re w=\pm W$ and odd
multiples of $\pi$ as horizontal sides, contained in
$|\Im w|\leq\nu H\widehat r$, the complete side estimates and
the homotopy $(1-s)+sA$ give the same zero count as $e^w-1$.
In particular the number of original zeros in $D(b,R)$ is greater than
\begin{equation}\label{bc:inverse-count}
 \nu H\widehat r/\pi-2>
 \nu(7/22)H\widehat r-2,
\end{equation}
provided the last expression is positive.  The loss two includes the
choice of the outermost odd heights.  When
$\nu H_0\widehat r>4\pi$, the top adjacent strip and the bottom
adjacent strip each contain one distinct simple original zero, with
\begin{equation}\label{bc:high-strips}
 |\Re w|<W,\qquad
 \Im w_+>\nu H\widehat r-4\pi,\qquad
 \Im w_-<-\nu H\widehat r+4\pi.
\end{equation}
These statements concern zeros of $P$ itself, through
\eqref{bc:inverse-location}.

\begin{lemma}[Localization of an edge normal]\label{bc:edge-normal}
Suppose $a$ is on a supporting edge of the original hull, the hull
contains a positive centered disk, and $|b-a|\leq e$.  Suppose
\eqref{bc:field-hypotheses}--\eqref{bc:inverse-gates-3} hold,
$\nu H_0\widehat r>4\pi$, and $|\Im S|<J_0$.
Set $s=J_0/H_0<1$ and
\begin{align}
 E_0&=W/H_0+s\nu\widehat r+e_R,\nonumber\\
 L_0&=(1-s^2)(\nu\widehat r-4\pi/H_0)-sW/H_0-e_R.
                                                        \label{bc:normal-errors}
\end{align}
If $L_0>e$, the outward normal $(\cos\theta,\sin\theta)$ of this
edge satisfies
\begin{equation}\label{bc:normal-bound}
 \cos\theta>0,\qquad |\tan\theta|<K:=\frac{E_0}{L_0-e}.
\end{equation}
In particular its distance from the origin is greater than
$a/\sqrt{1+K^2}$.
\end{lemma}

\begin{proof}
Write $\sigma=S/H$.  Since $\Re S>0$, one has
$|\Im\sigma|<s$ and $\Re\sigma>1-s^2$.  The two zeros supplied
by \eqref{bc:high-strips}, and $w/S=\overline\sigma w/H$, satisfy
\[
 \Re z_\pm>x-E_0,\qquad
 \Im z_+>y+L_0,\qquad \Im z_-<y-L_0.
\]
These bounds follow by applying \eqref{bc:inverse-location} on the
strips.
The positive support disk gives $a\cos\theta>0$.  Every original
zero obeys
\[
 \cos\theta(\Re z-a)+\sin\theta\Im z\leq0.
\]
Use $z_+$ when $\sin\theta\geq0$ and $z_-$ otherwise, and use
$x-a\geq0$, $|y|\leq e$.  This gives
$|\sin\theta|(L_0-e)<E_0\cos\theta$, proving the lemma.
\end{proof}

\begin{lemma}[Geometry of the whole edge and radial paths]
\label{bc:edge-geometry}
Suppose the edge through $a$ has distance at least $h>0$ from the
origin and its normal satisfies \eqref{bc:normal-bound}.  Set
\[
 c=\frac{h-K\sqrt{1-h^2}}{\sqrt{1+K^2}}.
\]
Every point $z$ of the edge has $\cos(\arg z)\geq c$ and
\begin{equation}\label{bc:whole-edge-gap}
 |z-\xi_j|>1-K-\sqrt{1-h^2}.
\end{equation}
Furthermore, suppose $a\geq a_1$, $0<\gamma<1$, and $0<r_0<h$.
If
\begin{equation}\label{bc:four-corners}
 1-v^2+2t(vc-\gamma)-t^2(1-\gamma^2)>0
 \quad(v=a_1,1;\ t=r_0,1),
\end{equation}
then the whole radial segment from $r_0u/|u|$ to each endpoint $u$
of the edge satisfies $|z-\xi_j|>\gamma|z|$.
\end{lemma}

\begin{proof}
If $h_e$ is the distance of the line and $\theta$ is its normal
angle, every point of its unit-disk chord satisfies
$|\arg z-\theta|\leq\arccos h_e\leq\arccos h$.
Combining this with $|\tan\theta|<K$ gives the bound $c$.
In normal and tangent coordinates the chord has tangent coordinate
in $[-\sqrt{1-h_e^2},\sqrt{1-h_e^2}]$, while the tangent
coordinate of $a$ has modulus $h_e|\tan\theta|$.  Therefore
$|z-a|\leq\sqrt{1-h^2}+K$, which proves
\eqref{bc:whole-edge-gap} by \eqref{bc:hole}.
For $z=t e^{i\arg u}$, the difference
\[
 (1-\gamma t)^2-|a-z|^2
 \geq1-a^2+2t(ac-\gamma)-t^2(1-\gamma^2)
\]
is bounded below by a polynomial concave separately in $a$ and $t$.
Its minimum on $[a_1,1]\times[r_0,1]$ is therefore at a corner.
Under \eqref{bc:four-corners}, $|a-z|<1-\gamma t$, and
\eqref{bc:hole} gives the claimed radial gap.
\end{proof}

\subsection{A potential estimate with an error proportional to the moment}

The following estimate retains a factor $Q$ in the interpolation error.

\begin{lemma}[Endpoint potential estimate]\label{bc:endpoint-potential}
Let $0<r_0<h$, and let $C_0,v_0,\beta,J,s_0,g,e_0$
be positive parameters.
Suppose an edge through a strict witness $a$ has distance at least $h$
from the origin, and all its points have critical distance greater than
$\gamma>0$.  Suppose, at its two endpoints $u$, the radial paths
from $r_0u/|u|$ to $u$ have critical distance greater than
$\gamma|z|$.  Assume
\[
 Q\leq q<1,\quad |\mu|\leq m<\min(h,r_0),\quad
 |P(0)|>p_f^n,\qquad p_f>.99.
\]
If $n\geq1000$, the comparisons
\eqref{bc:outer-gates}, \eqref{bc:endpoint-log-gates},
\eqref{bc:inner-gates}, and \eqref{bc:hessian-test} hold, and
$12+J/(1-Jq)+1/99<24$, then
\begin{equation}\label{bc:potential-24}
 -\log|P(0)|<24Q.
\end{equation}
\end{lemma}

\begin{proof}
Put $v=Q-|\mu|^2=N^{-1}\sum_j|\xi_j-\mu|^2\geq0$.
The exact centered resolvent identity is
\begin{equation}\label{bc:centered-resolvent}
 \frac1N\sum_j\frac1{z-\xi_j}-\frac1{z-\mu}
 =\frac1N\sum_j
 \frac{(\xi_j-\mu)^2}{(z-\mu)^2(z-\xi_j)}.
\end{equation}
On each outer radial path its modulus is at most $C_0v$ if
\begin{equation}\label{bc:outer-gates}
\begin{gathered}
 \frac1{(r_0-m)^2\gamma r_0}<C_0,\qquad
 1-\frac{m}{r_0(r_0-m)}>v_0,\qquad
 v_0\frac Nn-C_0q>\beta,\\
 \frac{C_0}{\beta^2(h-m)}<J,\qquad Jq<1.
\end{gathered}
\end{equation}
The second comparison bounds from below the inward logarithmic speed
of the centered radius $|z-\mu|$.  Indeed, for $z=t\omega$,
\[
 \Re\frac\omega{t\omega-\mu}
 \geq\frac1t-\frac{m}{t(t-m)}
 \geq1-\frac{m}{r_0(r_0-m)}.
\]
Let $q_P=P'/n$, $C_u=q_P(u)/(u-\mu)^N$, and
$w_u=r_0u/|u|$.  Integrate the logarithmic derivative in
\eqref{bc:centered-resolvent} inward from $u$ a distance $s$.
The centered monomial decays at least as $e^{-Nv_0s}$, and the
logarithm of its ratio to $q_P$ has modulus at most
$NC_0vs$.  Since $e^u-1\leq ue^u$ for $u\geq0$, the
absolute difference integral is at most
\[
 |P'(u)|\int_0^\infty nC_0vs e^{-\beta ns}\,ds
 =\frac{C_0v|P'(u)|}{\beta^2n}.
\]
Consequently
\begin{equation}\label{bc:outer-primitive}
\begin{gathered}
 \int_{w_u}^uP'(z)\,dz
 =C_u\bigl((u-\mu)^n-(w_u-\mu)^n\bigr)+E_{\rm out},\\
 |E_{\rm out}|\leq Jv|C_u(u-\mu)^n|.
\end{gathered}
\end{equation}

For the endpoint logarithm set $\zeta_j=(\xi_j-\mu)/(u-\mu)$.
Their mean is zero.  The parameter comparisons we use are
\begin{equation}\label{bc:endpoint-log-gates}
\begin{gathered}
 \frac{1+m}{h-m}<1.1,\quad
 \log\frac{1+m}{\gamma}<1.2,\quad
 \log\frac2{h-m}<1.2,\quad
 \frac{4(1.2+1.1)}{(h-m)^2}<12.
\end{gathered}
\end{equation}
They imply $|\zeta_j|<1.1$ and
$\gamma/(1+m)<|1-\zeta_j|<2/(h-m)$.
For $|\zeta|\leq1/2$, the complete logarithmic series gives
$|\log|1-\zeta|+\Re\zeta|\leq|\zeta|^2$.
For $|\zeta|>1/2$, the corresponding bound is
$4(1.2+1.1)|\zeta|^2$.  Averaging, and using the zero mean,
gives
\begin{equation}\label{bc:centered-log}
 |L(u)-\log|u-\mu||\leq12v.
\end{equation}

For the inner segment, write $\eta_j=\xi_j-\mu$ and, for fixed $z$,
consider
$f_z(t)=\prod_j(z-\mu-t\eta_j)$.  Its first derivative at $t=0$
vanishes.  For every product missing two labels, AM--GM gives
\[
 \prod_{j\ne i,k}|z-\mu-t\eta_j|
 \leq e\bigl(|z-\mu|^2+t^2v\bigr)^{(N-2)/2},
\]
because $(N/(N-2))^{(N-2)/2}\leq e$.  Also
$\sum_{i\ne k}|\eta_i\eta_k|\leq N^2v$.
Taylor's integral remainder therefore shows that
\[
 |q_P(z)-(z-\mu)^N|\leq2N^2vs_0^{N-2}
 \qquad(|z|\leq r_0),
\]
provided $(r_0+m)^2+q<s_0^2$.  The same argument applied to
$\prod_j(1-t\eta_j/(u-\mu))$ gives
\[
 |C_u|\leq e^{gn},\qquad
 |C_u-1|\leq2N^2ve^{gn}
\]
under $q/[2(h-m)^2]<g$ and $e/[2(h-m)^2]<2$.
If $(r_0+m)e^g<s_0$ and $e^gs_0^2<1$, it follows that
\[
 |q_P(z)-C_u(z-\mu)^N|\leq4N^2vs_0^{N-2}
 \qquad(|z|\leq r_0).
\]
Combining its integral from $0$ to $w_u$ with
\eqref{bc:outer-primitive} and $P(u)=0$ gives
\begin{equation}\label{bc:endpoint-identity}
\begin{gathered}
 0=P(0)+C_u(u-\mu)^n(1+\theta_u)+E_u,\quad
 |\theta_u|\leq Jv,\\
 |E_u|\leq4n^3vs_0^{n-3}+e^{gn}|\mu|^n.
\end{gathered}
\end{equation}
The condition $m<.006$ gives $|\mu|^n\leq Q(.006)^{n-2}$.
The inner conditions are
\begin{align}
 &(r_0+m)^2+q<s_0^2,\quad
   q/[2(h-m)^2]<g,\quad e<4(h-m)^2,\nonumber\\
 &(r_0+m)e^g<s_0,\quad e^gs_0^2<1,\quad
   4/s_0^3<6,\quad s_0/.99<e_0,\quad m<.006,\nonumber\\
 &e^g(.006)/.99<.01,\quad
   6(1000)^3e_0^{1000}+10^6(.01)^{1000}<.01,\quad
   e_0(1001/1000)^3<1.                            \label{bc:inner-gates}
\end{align}
They imply, for every $n\geq1000$,
\begin{equation}\label{bc:Q-error}
 |E_u/P(0)|\leq Q\bigl(6n^3e_0^n+10^6(.01)^n\bigr)\leq Q/100.
\end{equation}
Since $Jq<1$ and $|E_u/P(0)|<1$, we may take logarithms in
\eqref{bc:endpoint-identity}.  Using
$\log|C_u(u-\mu)^n|=NL(u)+\log|u-\mu|$, we obtain
\[
 n\left|L(u)-\frac1n\log|P(0)|\right|
 \leq\left(12+\frac J{1-Jq}+\frac1{99}\right)Q\leq24Q.
\]
The inequalities are strict when $Q>0$.

The assumed inequality \eqref{bc:hessian-test} implies $2h^2-1>0$.
The edge has unit tangent $v_e=ie^{i\theta}$ and support coordinates
$z=e^{i\theta}(h_e+it)$, with $h_e\geq h$ and
$h_e^2+t^2=|z|^2\leq1$.  Hence $h_e^2-t^2\geq2h_e^2-1>0$, and
\[
 -\Re\frac{v_e^2}{z^2}
 =\frac{h_e^2-t^2}{(h_e^2+t^2)^2}
 \geq h_e^2-t^2\geq2h_e^2-1\geq2h^2-1.
\]
The identity
\[
 (z-\xi)^{-2}=z^{-2}+2\xi z^{-3}
       +\frac{\xi^2(3z-2\xi)}{z^3(z-\xi)^2},
 \qquad 3z-2\xi=z+2(z-\xi),
\]
therefore bounds the second derivative of $L$ along the whole edge
from below by
\begin{equation}\label{bc:hessian-test}
 2h^2-1-\frac{2m}{h^3}
       -q\left(\frac1{h^2\gamma^2}+\frac2{h^3\gamma}\right)>0.
\end{equation}
Thus $L$ is convex on that edge.  Its two endpoint estimates imply
$nL(a)\leq\log|P(0)|+24Q$.  Since $L(a)>0$, this is
\eqref{bc:potential-24}.  The argument also covers $v=0$.
\end{proof}

\begin{lemma}[Annular cubic rigidity]\label{bc:cubic-rigidity}
Suppose all original zeros satisfy $\rho\leq|z_i|\leq1$, where
$0<\rho\leq1$, and all critical points belong to
$\{z:|z|\leq1,\ |z-1|\geq1\}$.  Suppose $C\geq1$, $0<q<1$,
$Q\leq q$, and $|P(0)|\geq e^{-CQ}$.  Then $Q=0$, and
$P(z)=z^n+c$ with $|c|=1$, provided
\begin{equation}\label{bc:rigidity-floor}
 n\geq\left\lceil\max\left\{
 \frac{256}{(1-q)^2},\ \frac{8Cq}{1-q},\
 C\left(\frac8\rho+\frac{16}{\rho^2}+\frac6{\rho^3}\right)+11
 \right\}\right\rceil.
\end{equation}
\end{lemma}

\begin{proof}
For $z=x+iy$ and $s=|z|^2$ the identity
\[
 \Re(-z-z^2-z^3/4)-s/4
 =(1+x)(3s/4-x-x^2)
\]
has nonnegative right side on the stated lens.  Indeed
$-1\leq x\leq1/2$, $s\geq x^2$, and $s\geq2x$; use the first
lower bound when $x\leq0$ and the second when $x\geq0$.
Writing $t_k=N^{-1}\sum_j\xi_j^k$, we deduce
\begin{equation}\label{bc:cubic-positive}
 Q\leq4|t_1|+4|t_2|+|t_3|.
\end{equation}

Assume $Q>0$ and put $a_j=P^{(j)}(0)/P(0)$, $j=1,2,3$.
Differentiate the product for $P'$ and apply AM--GM to every remaining
subproduct.  Since $(N/(N-j+1))^{(N-j+1)/2}\leq e^{(j-1)/2}<3$,
\[
 |a_j|\leq3e^{Cq}n^3Q^{(n-3)/2}.
\]
Set $\delta=1-q$.  The first two conditions in
\eqref{bc:rigidity-floor} imply
\[
 3e^{Cq}n^4q^{(n-5)/2}<1.
\]
$\log n\leq\sqrt n$ and $\log q\leq-\delta$ give
\begin{align*}
 \log3+Cq+4\log n+\tfrac{n-5}2\log q
 &\leq\log3+Cq-\delta n/4+5\delta/2\\
 &\leq\log3-\delta n/8+5\delta/2\\
 &\leq\log3-32/\delta+5\delta/2<0.
\end{align*}
It follows that $|a_j|\leq Q/n$.

Let $R_k=\sum_i z_i^{-k}$ and $s_k=n^{-1}\sum_i z_i^k$.  The exact
identities
\[
 R_1=-a_1,\quad R_2=a_1^2-a_2,\quad
 R_3=-(a_3-3a_1a_2+2a_1^3)/2
\]
give $|R_k|\leq2Q/n$ for $k=1,2,3$.  For an original modulus
$s\in[\rho,1]$, integrate the derivative of $e^{kt}-e^{-kt}$ from
$0$ to $t=-\log s$ to obtain
\[
 |z_i^k-\overline{z_i^{-k}}|
 =s^{-k}-s^k\leq2k\rho^{-k}(-\log s).
\]
Sum over all original zeros.  The radial budget
$\sum_i-\log|z_i|\leq CQ$ implies
\begin{equation}\label{bc:power-moment-small}
 |s_k|\leq A_kQ/n,\qquad A_k=2kC\rho^{-k}+1,\quad k=1,2,3.
\end{equation}
Here $n\geq256$ justifies replacing the term $2/n$ by one.

Newton's identities for $P$ and $P'/n$ give
\[
 t_1=s_1,\qquad
 t_2=\frac{(n-2)s_2+s_1^2}{n-1},\qquad
 t_3=\frac{(n-3)s_3+3s_1s_2-s_1^3}{n-1}.
\]
Using $|s_1|\leq1$, \eqref{bc:cubic-positive}, and
\eqref{bc:power-moment-small}, we obtain
\[
 Q\leq\left(B_1+9+\frac{E_1}{n-1}\right)\frac Qn,
\]
where
\[
 B_1=C(8/\rho+16/\rho^2+6/\rho^3),\qquad
 E_1=10C/\rho+12C/\rho^2+8.
\]
But $B_1+10-E_1=C(-2/\rho+4/\rho^2+6/\rho^3)+2>0$.
Since $n\geq B_1+11$, the coefficient in parentheses is less than
$B_1+10<n$, a contradiction.  Thus $Q=0$.
Now $P'=nz^{n-1}$, and the product hypothesis and the unit-disk roots
force the constant term of $P$ to have modulus one.
\end{proof}

\begin{proposition}[Boundary and small displacement ranges]\label{bc:boundary}
For every $n\geq6251$, a hull point satisfying \eqref{bc:hole} is
impossible in each of the following cases:
\begin{equation}\label{bc:boundary-scopes}
\begin{array}{c|c|c}
 &a&r\\ \hline
 \mathrm B&[.975,1]&\text{arbitrary}\\
 \mathrm A&[.85,1]&[0,.02]\\
 \mathrm {G1}&[.96,1]&[0,.04]\\
 \mathrm {G2}&[.965,1]&[0,.045].
\end{array}
\end{equation}
The latter three conclusions hold for any choice of actual rightmost
zero satisfying the indicated displacement bound.
\end{proposition}

\begin{proof}
We first treat a boundary witness $a<1$ which is not an original
zero.  The normalization and moment proof above gives
Table~\ref{bc:refined-table}.  Apply Lemma~\ref{bc:poisson-support}
using the weak support radii in Table~\ref{bc:geometry-table}.
For row A, repeat the support test and the actual-parameter moment
inequality in the following order:
\begin{equation}\label{bc:feedback}
\begin{array}{c|c|c}
 \text{lower }a&\text{new upper }Q&\text{next support radius}\\ \hline
 .9305&.13831&.9454\\
 .9454&.11036&.9478\\
 .9478&.10582&.9481\\
 .9481&.10525&\text{--}.
\end{array}
\end{equation}
Each new moment bound is obtained from \eqref{bc:moment-Q} by
increasing only the lower bound for its explicit $a^2$ term.
The error $\mathcal E(.02;.85,N)$ is retained throughout, and the
caps for $B,|\mu|$ remain unchanged.
Put $a_1=.9481$, $q=.10525$ in row A;
in the other three rows put $a_1=a_0$, $q=q_0$.

At $b$ the finite resolvent expansion
\[
 \frac SN=\frac1b+\frac{\mu}{b^2}
                  +\frac1N\sum_j\frac{\xi_j^2}{b^2(b-\xi_j)}
\]
gives
\begin{equation}\label{bc:moment-field}
 H>Nc_f,\quad E_b/H<V,\quad d^{-1}<M,\qquad
 c_f=1-\frac m{a_1}-\frac q{a_1(1-e)},\quad
 V=\frac1{c_f(1-e)^2},\quad M=\frac1{1-e}.
\end{equation}
The listed parameters give $c_f>0$.  Since
$1-m/|b|-q/(|b|d)>c_f$ and $|b|\leq1$,
\[
 H\geq\frac N{|b|}\left(1-\frac m{|b|}-\frac q{|b|d}\right)>Nc_f.
\]
Use $H_0=N_0c_f$, $\nu=.99$, $W=4$ and the radii $R$ in
Table~\ref{bc:geometry-table}.  All comparisons
\eqref{bc:inverse-gates-1}--\eqref{bc:inverse-gates-3} hold, as does
$\nu H_0\widehat r>4\pi$.  Equation \eqref{bc:fine-phase} supplies
$J_0=8/a_0$, with the original $a_0$ also in row A.
Lemma~\ref{bc:edge-normal} consequently gives $K$.

For B, G1, and G2 the direct comparisons
$a_0^2>h^2(1+K^2)$ show that this particular edge has distance
greater than the listed $h$.  Row A uses its support radius $.9481$.
All four rows satisfy
$1-K-\sqrt{1-h^2}>\gamma$ and the four corner inequalities
\eqref{bc:four-corners} with $r_0=.9$.
Lemma~\ref{bc:edge-geometry} therefore supplies the gaps on the
whole edge and the endpoint radial paths.

The endpoint parameters in Table~\ref{bc:endpoint-table} are positive,
and $r_0=.9<h$ in all four rows.
All the hypotheses of Lemma~\ref{bc:endpoint-potential} hold.
In particular the Hessian bounds and endpoint coefficients in
Table~\ref{bc:boundary-margins} are strictly positive and strictly
less than $24$, respectively.  We obtain
$-\log|P(0)|<24Q$.
The additional annulus radii $\rho_f$ in the table follow from
\[
 \rho_f^2+2\rho_fm+q<s_f^2,\qquad
 \rho_f<s_f,\qquad
 s_f/p_f<s_*/p,
\]
and the already proved whole-degree bound $n(s_*/p)^n<1$.
Thus every original zero satisfies $|z_i|>\rho_f$.

Every critical point belongs to the lens in
Lemma~\ref{bc:cubic-rigidity}, since
\[
 |\xi-1|^2-|\xi-a|^2=(1-a)(1+a-2\Re\xi)>0.
\]
The three bounds in \eqref{bc:rigidity-floor}, with $C=24$,
$\rho=\rho_f$, and the final $q$, are all less than $900$.
Lemma~\ref{bc:cubic-rigidity} gives $Q=0$.  All critical points
then equal zero, contradicting \eqref{bc:hole} for $a<1$.

The weak support test already excludes a hull of dimension less than
two.  A boundary vertex is an original zero and was excluded by
the energy inequality; every other boundary point lies on a genuine
edge as used above.  If $a=1$, strict convexity of $\overline\D$
makes $1$ an original zero, again impossible.
Finally, for an interior witness move outward on its real ray to the
boundary point $A$.  Since $a\leq A\leq x$,
\[
 |A-\xi|^2-|a-\xi|^2=(A-a)(A+a-2\Re\xi)\geq0,\qquad
 |b-A|\leq|b-a|.
\]
The strict hole and each applicable displacement hypothesis persist.
The boundary result therefore proves all of \eqref{bc:boundary-scopes}.
\end{proof}

\begin{table}[ht]
\centering\small
\caption{Support and geometric parameters for the four exclusion rows.}
\label{bc:geometry-table}
\begin{tabular}{@{}crrrrrrr@{}}\toprule
Row&weak $h$&final $h$&$R$&$\gamma$&$\rho_f$&$s_f$&final $q$\\\midrule
B & 0.9239 & 0.95027 & 0.113 & 0.455 & 0.955 & 0.993 & 0.06001\\
A & 0.9305 & 0.9481 & 0.061 & 0.583 & 0.934 & 0.993 & .10525\\
G1 & 0.9351 & 0.94709 & 0.092 & 0.51 & 0.943 & 0.993 & 0.08531\\
G2 & 0.9308 & 0.94828 & 0.1 & 0.49 & 0.947 & 0.993 & 0.07689\\
\bottomrule\end{tabular}
\end{table}
\begin{table}[ht]
\centering\small
\caption{Endpoint parameters, with $r_0=.9$ in every row.}
\label{bc:endpoint-table}
\begin{tabular}{@{}crrrrrrr@{}}\toprule
Row&$C_0$&$v_0$&$\beta$&$J$&$s_0$&$g$&$e_0$\\\midrule
B & 3.056 & 0.9932 & 0.8086 & 4.952 & 0.93797 & 0.03371 & 0.947455\\
A & 2.369 & 0.9973 & 0.7468 & 4.495 & 0.95869 & 0.0589 & 0.968384\\
G1 & 2.716 & 0.9955 & 0.7626 & 4.954 & 0.94956 & 0.048 & 0.959162\\
G2 & 2.83 & 0.9947 & 0.7759 & 4.984 & 0.94572 & 0.04323 & 0.955283\\
\bottomrule\end{tabular}
\end{table}
\begin{table}[ht]
\centering\small
\caption{Outward numerical bounds for the boundary geometry.
The comparisons themselves use rational intervals.}
\label{bc:boundary-margins}
\begin{tabular}{@{}crrr@{}}\toprule
Row&upper bound for $K$&lower Hessian bound&upper endpoint coefficient\\\midrule
B&.229122900&.165248328&19.055896748\\
A&.092635534&.024898554&20.541111328\\
G1&.164970033&.026449841&20.590324175\\
G2&.187982215&.064807027&20.090774931\\\bottomrule
\end{tabular}
\end{table}

The first support comparison and each feedback comparison
check both \eqref{bc:support-test} at $n_0$ and its slope
\[
 \Lambda(1-h)^2-2B_0>0.
\]
The four inverse applications check
\eqref{bc:inverse-gates-1}--\eqref{bc:inverse-gates-3} for both signs
of the strip, as well as $s<1$, $\nu\widehat r-4\pi/H_0>0$,
and $L_0>e$.  The geometric comparisons are $h>.9$,
$1-K-\sqrt{1-h^2}>\gamma$, $c>0$, and the four instances of
\eqref{bc:four-corners}.  The endpoint checks are exactly
\eqref{bc:outer-gates}, \eqref{bc:endpoint-log-gates},
\eqref{bc:inner-gates}, the stricter lower bound $.01$ in
\eqref{bc:hessian-test}, and
$12+J/(1-Jq)+1/99<24$.
The final annulus and the three inequalities in
\eqref{bc:rigidity-floor} are checked as in the proof.
These tests, together with \eqref{bc:pre-gates-1}--\eqref{bc:pre-gates-9},
form the boundary certificate.

As $n$ increases, the fixed moment caps remain valid, while the field
$Nc_f$ increases.  Thus the inverse corner inequality improves and
$\epsilon$ decreases.  In \eqref{bc:normal-errors}, $E_0$ decreases
and $L_0$ increases, so $K$ decreases.  The geometric inequalities
therefore persist.  The outer rate $v_0N/n-C_0q$ increases, and the
inner errors have the explicit decreasing consecutive-degree ratios
in \eqref{bc:inner-gates}.  This proves the boundary conclusions on
the whole infinite degree range, including all stated endpoints.

\subsection{A direct entrance for the compact range}

\begin{lemma}[Crescent entrance]\label{bc:crescent}
Under \eqref{bc:hole}, if $a\geq.1$ and $n\geq6251$, then
\begin{equation}\label{bc:crescent-D}
 D:=Hd>\frac{5N}{14}\frac{c}{1-c+c^2}>117.17,
 \qquad c=x/2.
\end{equation}
Thus Propositions~\ref{mf:D100} and \ref{lt:D70} apply.
\end{lemma}

\begin{proof}
A closest point to $b=x+iy$, $x\geq a$, $y\geq0$, in the closed
crescent $\{|\xi|\leq1,\ |\xi-a|\geq1\}$ is
$a/2+i\sqrt{1-a^2/4}$.  Indeed a minimizer must be on one of its
two circular boundary arcs.  On $|\xi|=1$, maximizing
$\Re(\overline b\xi)$ over $\Re\xi\leq a/2$ gives the stated
endpoint, because the direction of $b$ lies to its right.
On $|\xi-a|=1$, the same argument maximizes
$\Re((\overline b-a)(\xi-a))$ over
$\Re(\xi-a)\leq-a/2$; $b-a$ has nonnegative real and imaginary
parts, giving the same endpoint.  Since the strict hole excludes
that endpoint,
\[
 d^2>1+x^2+y^2-ax-2y\sqrt{1-a^2/4}.
\]
The right side decreases in $a$ for $0<a\leq x\leq1$, since its
derivative is $-x+ay/[2\sqrt{1-a^2/4}]<0$.
Set $c=x/2$, $v=\sqrt{1-c^2}-y\geq0$, and
$G=\sqrt{c^2+v^2}$.  Comparing at $a=x$ gives $d>G$.
Also, if $p=1-|b|^2$, then
\[
 p=-3c^2+2v\sqrt{1-c^2}-v^2\leq2v-3c^2-v^2.
\]
The exact radial identity, obtained by summing
$2\Re(b/(b-\xi))+(1-|b|^2)/|b-\xi|^2
=1+(1-|\xi|^2)/|b-\xi|^2$, is
\begin{equation}\label{bc:radial-identity}
 N=2\Re(bS)+(1-|b|^2)E_b
       -\sum_j\frac{1-|\xi_j|^2}{|b-\xi_j|^2}.
\end{equation}
Furthermore $\Re S>cE_b$, because $a\leq x$.
Using $|b|\leq1-p/2$ in \eqref{bc:radial-identity}, we obtain
\[
 \frac NH\leq2+p(c^{-1}-1),\qquad
 Hd>\frac{NcG}{2c+(1-c)p}.
\]
The second inequality is strict because $d>G$.
By Cauchy--Schwarz and the identity
$(3c^2+v^2)^2-8c^2(c^2+v^2)=(v^2-c^2)^2$,
\begin{align*}
 \frac{2c+(1-c)p}{G}
 &\leq\frac{2c+2(1-c)v}{G}
       -(1-c)\frac{3c^2+v^2}{G}\\
 &\leq2\sqrt{2-2c+c^2}-2\sqrt2\,c(1-c).
\end{align*}
For $.05\leq c\leq1/2$, $\sqrt{2-2c+c^2}<7/5<\sqrt2$.
This proves the first inequality in \eqref{bc:crescent-D}.
The function $c/(1-c+c^2)$ increases on that interval.
Its value at $c=.05$, with $N=6250$, proves the second.
\end{proof}

\begin{proposition}[The bridge interval]\label{bc:bridge}
For every $n\geq6251$, no hull witness satisfies
\eqref{bc:hole} with $.1\leq a\leq.11$.
\end{proposition}

\begin{proof}
Lemma~\ref{bc:crescent} gives $D>100$.  Throughout the argument,
$C$ and $\gamma$ are the quantities in \eqref{bc:adaptive-def}.
The logarithmic form of Lemma~\ref{bc:blaschke} is
\begin{equation}\label{bc:log-bridge}
 -\log\mathcal B_P(\gamma)+CE_b/2<\log(8/C).
\end{equation}

\emph{The range $|b|\leq.2$.}
Here $d>.8$ and Proposition~\ref{mf:D100} gives
$E_bd/H<9/4$.  In particular
$V=E_b/H<9/(4\cdot.8)$ and
$H>N(.8)/(9/4)$, using $E_b\geq N$.
Apply \eqref{bc:inverse-gates-1}--\eqref{bc:inverse-gates-3}
with $R=.1$, $\nu=.99$, $W=6$, $e_W=1/400$,
$M=1/.8$, and $H_0=N_0(.8)/(9/4)$.
The image radius is $\widehat r=379/4480$, and the count
\eqref{bc:inverse-count} is greater than $1831/32$ at the floor.
Since $|\gamma|\leq|b|$, each counted zero has factor at most
\[
 \frac{R+.11}{1-(.2+R)(.2)}<\frac14.
\]
As $e<4$, its logarithmic weight exceeds one.
The count divided by $n$ increases, and exceeds $1/120$ for all
$n\geq6251$.

The actual-zero variance inequality \eqref{ar:variance-unit}
gives $d^2\leq1-(1-|b|^2)/n$, and therefore $d<1-.47/n$.
Since $d+r>1$, we have $r>.4/n$.
Since $a\geq.1$ and $u=\min(r^2,1/2)>(.4/n)^2$,
the monotonicity of $u(1-u)$ on $[0,1/2]$ gives
\[
 C\geq a^2u(1-u)
 >(.1)^2\left(\frac{.4}{n}\right)^2
       \left(1-\left(\frac{.4}{n}\right)^2\right)
 >\frac1{1000n^2}.
\]
Hence
\[
 \log(8/C)<10+2\log n<n/120
       <-\log\mathcal B_P(\gamma),
\]
a contradiction to \eqref{bc:log-bridge}.
For the middle inequality, at $n=6251$ one has
$\log n<9$ and $n/120>28$; its difference increases since
$1/120>2/n$.

\emph{The range $.2<|b|\leq.8$.}
Now $r>.09$, so
$C>C_m=(.1)^2(.09)^2(1-(.09)^2)$.
The field conclusion of Proposition~\ref{lt:D70}, with $D>100$
from Lemma~\ref{bc:crescent}, gives
\[
 H>\frac{N-(1899/20)/(.2)}{2(1+.11)}.
\]
Use $d>.2$, $V=9/(4\cdot.2)$, $M=1/.2$,
$R=.04$, and the same $\nu,W,e_W$.
The image radius is $.0295$, the count at the floor exceeds
$26254139/1184000$, and every counted factor is at most
\[
 \frac{.04+.11}{1-(.8+.04)(.8)}<.46.
\]
The comparison $e^{.75}<1/.46$ gives a product budget exceeding
$16.63>12$, while $e^{12}>8/C_m$.
This again contradicts \eqref{bc:log-bridge}; its count increases
with $N$.

\emph{The range $.8<|b|\leq.94$.}
Here $r>.69$, so $u\geq.4761$, $1/2\leq t\leq.5239$, and
\[
 C>C_u=(.1)^2(.4761)(.5239)=24942879/10^{10}.
\]
Retain only the factor of the original zero $b$ itself.
Its numerator is $ta\leq.057629$, while the real part of its
denominator is
$1-|b|^2+tax\geq1-.94^2+.5(.1)^2=.1214$.
Thus that factor is at most $57629/121400<.475$, and its logarithmic
weight exceeds $.74$.  Since
\[
 .74+6250C_u/2>8.08,\qquad e^{8.08}>8/C_u,
\]
the product factor and the energy term for the same $C$ contradict
\eqref{bc:log-bridge}.

\emph{The range $|b|>.94$.}
Now $r>.83$, hence $u=t=1/2$.
Put $a_*=.03$, $x_* =.1$, and $s=3a_*/4$.
Concavity of $a(x-3a/4)$ on $[a_*,x]$, and $x>3a_*$, give
\[
 C=a(x-3a/4)\geq a_*(x-s).
\]
In fact the value at $a=x$ exceeds the value at $a=a_*$ because
$(x-a_*)(x-3a_*)>0$.
The comparison $e^{8.15}>8/[a_*(x_*-s)]$ and
\eqref{bc:log-bridge} imply
$E_b/N<k/(x-s)$, where $k=2(8.15)/(Na_*)$.
Also $H^2\leq NE_b$ by Cauchy--Schwarz.
The phase conclusion of Proposition~\ref{lt:D70} is
$|\Im S|<(21/4)E_b/H$; the half-plane bound
$E_b/H<2/x$ follows directly from $\Re S>(x/2)E_b$.
Use these estimates in the exact identity \eqref{bc:radial-identity}.
Writing $p_{\max}=1-.94^2$, we find the necessary inequality
\begin{equation}\label{bc:high-necessary}
 1<2x\sqrt{\frac{k}{x-s}}
       +\frac{p_{\max}k}{x-s}
       +\frac{4(21/4)}{Nx}.
\end{equation}
Under $v=\sqrt{x-s}$ the right side is a positive linear combination
of $v+s/v$, $v^{-2}$ and $(v^2+s)^{-1}$.
Each is convex throughout the interval in question: for the last
term this follows from $3v^2>s$.
Its maximum is therefore at $x=x_*$ or $x=1$.
At each of those endpoints, with $N=N_0$, the exact comparisons
\[
 R_x:=1-\frac{p_{\max}k}{x-s}-\frac{21}{Nx}>0,\qquad
 R_x^2>\frac{4x^2k}{x-s}
\]
make the right side of \eqref{bc:high-necessary} less than one.
Every term decreases with $N$, which proves the contradiction in
the whole high-modulus range.  The four cases include their common
endpoints and exhaust every possible $b$.
\end{proof}

\subsection{Nested original-zero counts in the middle range}

\begin{proposition}[The compact range]\label{bc:compact}
For every $n\geq6251$, no hull witness satisfies
\eqref{bc:hole} with $.1\leq a\leq.975$.
\end{proposition}

\begin{proof}
We prove the part $.11\leq a\leq.975$ by the closed cover in
Table~\ref{bc:compact-cover}; Proposition~\ref{bc:bridge} supplies
the remaining interval.  A row specifies $A\leq a\leq B$ and
$r_-(n)\leq r\leq r_+$.
All numbers in the tables are exact rationals.

We first construct a lower bound $H(n)=\alpha(n-1)-\beta$ for $H$,
together with constants $V,M$ as in \eqref{bc:field-hypotheses}.
There are two constructions.

For a trace row put
\[
 M=\frac1{1-r_+},\qquad c=\frac{A}{2(1+r_+)^2}.
\]
Lemma~\ref{bc:crescent} gives $D>100$.
Since $1-|b|^2\leq2d$, \eqref{bc:radial-identity} and
$E_bd/H<9/4$ give $H>2N/13$ and
$D>2N(1-r_+)/13$.
Set
\begin{equation}\label{bc:Dfloor}
 D_f=\frac{2N_0(1-r_+)}{13}-10^{-6}.
\end{equation}
For the six trace rows $D_f\geq201$.
The Fine signed-moment inequality \eqref{lt:fine-curvature}, derived
from Lemma~\ref{lt:signed-moment},
on its domain $E_b/H^2\leq7/625$, $(Hd)^{-1}\leq1/200$,
applies here: $E_b/H^2<(9/4)/201<7/625$.
For its ratio consequence, set
$v_j=[H(b-\xi_j)]^{-1}$, $q_b=E_b/H^2$,
$\mu_b=(Hd)^{-1}$, and $y_b=E_bd/H=q_b/\mu_b$.
Then
\[
 q_b-2\mu_b
 \leq-\Re\sum_jv_j^2
 <(9/4)q_b^2+13q_b\mu_b,
\]
because $0\leq\Re v_j\leq\mu_b$ and
$\sum_j\Re v_j\leq1$.
Consequently
\begin{equation}\label{bc:ratio-quadratic}
 (9/4)y_b^2-(D-13)y_b+2D>0.
\end{equation}
The quadratic is strictly decreasing on $[2,9/4]$ for $D\geq201$.
The values $\kappa$ in Table~\ref{bc:trace-ratios} satisfy
\begin{equation}\label{bc:kappa-gate}
 2<\kappa<9/4,\qquad
 (9/4)\kappa^2-(D_f-13)\kappa+2D_f<0.
\end{equation}
As its value at $\kappa>2$ decreases with $D$, and $y_b<9/4$ is
already known, \eqref{bc:ratio-quadratic} implies $y_b<\kappa$.
We may therefore take
\[
 V=\min(2/A,\kappa/(1-r_+)),\qquad
 H_{\rm init}=\max(N_0A/2,2N_0/13).
\]
The comparisons $V/H_{\rm init}<7/625$ and
$M/H_{\rm init}<1/270$ justify the signed estimate
\[
 \Re\sum_j(b-\xi_j)^{-2}>-(9/4)V^2-13VM.
\]
Each real reciprocal coordinate
$X_j=\Re(b-\xi_j)^{-1}$ belongs to $(c,M)$:
\[
 2(a-\Re\xi)-a|a-\xi|^2
 =a(1-|\xi|^2)+(1-a^2)(a-2\Re\xi)>0,
\]
and $|b-\xi|<(1+r_+)|a-\xi|$.
Using $X_j^2\leq(c+M)X_j-cM$ and $E_b\geq N$ in
$E_b=2\sum_jX_j^2-\Re\sum_j(b-\xi_j)^{-2}$ gives
\begin{equation}\label{bc:trace-field}
 H>\alpha N-\beta,\quad
 \alpha=\frac{1+2cM}{2(c+M)},\qquad
 \beta=\frac{(9/4)V^2+13VM}{2(c+M)}.
\end{equation}

For a moment row take $(\ell,w,\kappa)$ from M0, M1, or M2 as
indicated in Table~\ref{bc:compact-cover}.  Define, at the fixed floor,
\begin{align}
 e_r&=\frac{8r_+}{AN_0}
            +\frac{r_+^2}{2(1-r_+)(1-2r_+)},\nonumber\\
 q_m&=1-A^2+w(\ell/N_0+e_r)+2/n_0^2,\qquad
 m_m=\kappa(\ell/n_0+e_r)+1/n_0^2.                \label{bc:compact-moments}
\end{align}
These are upper bounds for $Q,|\mu|$ throughout the row
and for every larger degree, by Proposition~\ref{bc:moment-proposition}.
Set
\begin{equation}\label{bc:compact-moment-field}
 \alpha=1-\frac{m_m}A-\frac{q_m}{A(1-r_+)},\quad
 \beta=0,\quad V=\frac1{\alpha(1-r_+)^2},\quad
 M=\frac1{1-r_+}.
\end{equation}
The finite comparisons give $\alpha>0$; the resolvent estimate
\eqref{bc:moment-field} gives $H>\alpha N$.

For each row and each indicated integer $j$, use $R_j=j/1000$ and
the parameters $\nu_j,W_j$ specified in
Tables~\ref{bc:nu-table} and \ref{bc:W-table}.
All intermediate integers between the displayed endpoints are included.
Use \eqref{bc:inverse-def-1}--\eqref{bc:inverse-def-2} with this row's
$V,M$, and set $H_0=\alpha N_0-\beta$.
Every one of the comparisons
\eqref{bc:inverse-gates-1}--\eqref{bc:inverse-gates-3} holds, together
with $\nu_jH_0\widehat r_j>22/7$.
Equation \eqref{bc:inverse-count} consequently gives, for the actual
number $\mathcal N_j$ of original zeros in $D(b,R_j)$,
\begin{equation}\label{bc:nested-count}
 \mathcal N_j>k_j(n):=
      \nu_j(7/22)\widehat r_j\bigl(\alpha(n-1)-\beta\bigr)-2>0.
\end{equation}

We estimate their Blaschke factors at $\gamma$ from
\eqref{bc:adaptive-def}.  In every middle row
$r^2\leq r_+^2<1/2$, so $t=1-r^2$.
For a counted zero $z=b+v$, rightmostness gives $\Re v\leq0$.
It follows that
\[
 |z-\gamma|^2=|v+ta|^2\leq R_j^2+B^2.
\]
Substitution of $C=ta[2x-(1+t)a]$ gives the denominator identity
\begin{equation}\label{bc:correlated-denominator}
 \Re(\overline b\gamma)
 =\frac{C(1+r^2)}{2(1-r^2)}
        +r^2+\frac{a^2r^2(1-r^2)}2.
\end{equation}
Also
$\Re(\overline v\gamma)=(x-ta)\Re v+y\Im v\leq r_+R_j$;
here $x-ta\geq0$.  Thus, with
\begin{equation}\label{bc:Delta-def}
 \Delta_j=
 \frac{C_*(1+r_+^2)}{2(1-r_+^2)}
       +(1+B^2/2)r_+^2+R_jr_+,
\end{equation}
we have $\Re(\overline z\gamma)<\Delta_j$.
The row comparisons give $0<\Delta_j<1$ and hence
\begin{equation}\label{bc:factor-bound}
 \left|\frac{\gamma-z}{1-\overline z\gamma}\right|
 \leq\frac{\sqrt{B^2+R_j^2}}{1-\Delta_j}\leq\eta_j<1.
\end{equation}
Here $\eta_j$ is the smallest multiple of $10^{-12}$ whose square is
at least $(B^2+R_j^2)/(1-\Delta_j)^2$.

Put $t_j=(1-\eta_j)/(1+\eta_j)$ and define the rational weights by
\begin{equation}\label{bc:weights}
 w_j=10^{-12}\left\lfloor10^{12}
       \left(2\sum_{k=0}^{23}\frac{t_j^{2k+1}}{2k+1}\right)
       \right\rfloor.
\end{equation}
All omitted terms of the logarithmic series are positive, so
$0<w_j<-\log\eta_j$.
For each row the $w_j$ decrease and the $k_j(n_0)$ increase with $j$.
The physical disks are nested, so finite Abel summation gives
\begin{align}
 -\log\mathcal B_P(\gamma)
 &>w_J\mathcal N_J+
         \sum_{j<J}(w_j-w_{j+1})\mathcal N_j\nonumber\\
 &>w_Jk_J(n)+
         \sum_{j<J}(w_j-w_{j+1})k_j(n)
   =\sum_jw_j\bigl(k_j(n)-k_{j-1}(n)\bigr)
   =:\mathcal A(n),                               \label{bc:Abel}
\end{align}
where $k_0=0$ and the indices are understood in their successive
order within a row.  The first inequality assigns to each counted
zero the largest available weight.  The second uses nonnegative
coefficients and lower bounds on nested disk counts.

Set
\[
 C_-(n)=A^2r_-(n)^2(1-r_-(n)^2).
\]
Because $r^2<1/2$, $C\geq C_-(n)$.
For a fixed radial lower endpoint retain the energy bonus
$(n-1)C_-/2$ in \eqref{bc:log-bridge}; for a shrinking lower
endpoint $r_-(n)=c/n$ discard that positive term.
The targets $T_0$ in Table~\ref{bc:compact-margins} satisfy
\begin{equation}\label{bc:compact-final-gates}
 e^{T_0}>\frac8{C_-(n_0)},\qquad
 \mathcal A(n_0)+\text{bonus}>T_0.
\end{equation}
Put $\lambda_j=\nu_j(7/22)\widehat r_j$ and
\[
 W_*=\sum_jw_j(\lambda_j-\lambda_{j-1}),\qquad\lambda_0=0.
\]
The exact telescoping identity is
\[
 \mathcal A(n)=W_*\bigl(\alpha(n-1)-\beta\bigr)-2w_{\rm first}.
\]
The comparisons verify that $\mathcal A(n)+\text{bonus}=A_*n-B_*$
with $A_*,B_*>0$.
For fixed $r_-$ the left side increases, while $\log(8/C_-)$ is
constant, so \eqref{bc:compact-final-gates} contradicts
\eqref{bc:log-bridge} for every $n\geq n_0$.
For $r_-(n)=c/n$,
$C_-(n)\geq C_-(n_0)(n_0/n)^2$, while $\mathcal A(n)/n$
increases.  As $T_0>2$ in these four rows,
\[
 \log\frac8{C_-(n)}
 <T_0+2\log(n/n_0)
 \leq T_0(n/n_0)<\mathcal A(n),
\]
again a contradiction on the entire degree range.
The field entrances and every inverse contour condition also improve
with $n$, since $H(n)$ increases and the other contour constants are
fixed.

It remains to verify that the closed cells just excluded leave no
gaps.  For each interval in $a$, its final far radius $r_f$ is listed
below:
\begin{equation}\label{bc:far-table}
\begin{array}{c|rrrrrrrrrr}
 A&.11&.2&.45&.75&.85&.9&.93&.95&.96&.965\\ \hline
 r_f&.65&.27&.115&.069&.06&.057&.055&.054&.053&.053.
\end{array}
\end{equation}
Each satisfies $r_f^2<1/2$ and
$A^2r_f^2(1-r_f^2)>C_*$.  Lemma~\ref{bc:blaschke} therefore
excludes every $r\geq r_f$, including $r^2\geq1/2$.
For the four shrinking lower endpoints, let
$\overline B=B+.001$.  If $r\leq c/n$, then $|b|<\overline B$,
and \eqref{ar:variance-unit} gives
\[
 d^2<1-(1-\overline B^2)/n,\qquad
 d<1-\frac{1-\overline B^2}{2n}<1-c/n.
\]
Here the four exact comparisons are $B+c/n_0<\overline B$ and
$c<(1-\overline B^2)/2$.
This contradicts $d+r>1$.
For $a\geq.85$ the interval $r\leq.02$ is excluded by row A of
Proposition~\ref{bc:boundary}.  For $a\geq.96,r\leq.04$, and
$a\geq.965,r\leq.045$, use rows G1 and G2, respectively.
All witness and radial endpoints are shared by closed exclusions.
Together with Proposition~\ref{bc:bridge}, this proves
Proposition~\ref{bc:compact}.
\end{proof}

\begin{table}[ht]
\centering\small
\caption{The thirteen closed middle cells.  Every integer in a displayed
$j$ interval is used, with $R_j=j/1000$.}
\label{bc:compact-cover}
\begin{tabular}{@{}crrrrrrl@{}}\toprule
Row&$A$&$B$&$r_-(n)$&$r_+$&$j$ interval&number&field\\\midrule
Ia & 0.11 & 0.2 & $0.45/n$ & 0.1 & 5--266 & 262 & Trace\\
Ib & 0.11 & 0.2 & 0.1 & 0.65 & 6--91 & 86 & Trace\\
II & 0.2 & 0.45 & $0.35/n$ & 0.27 & 5--212 & 208 & Trace\\
III & 0.45 & 0.75 & $0.2/n$ & 0.115 & 5--262 & 258 & Trace\\
IV & 0.75 & 0.85 & $0.12/n$ & 0.069 & 5--278 & 274 & Trace\\
V & 0.85 & 0.9 & 0.02 & 0.06 & 5--281 & 277 & Trace\\
VI & 0.9 & 0.93 & 0.02 & 0.057 & 5--299 & 295 & M0\\
VII & 0.93 & 0.95 & 0.02 & 0.055 & 5--244 & 240 & M1\\
VIIIa1 & 0.95 & 0.96 & 0.02 & 0.025 & 5--248 & 244 & M2\\
VIIIa2 & 0.95 & 0.96 & 0.025 & 0.035 & 5--236 & 232 & M2\\
VIIIb & 0.95 & 0.96 & 0.035 & 0.054 & 5--211 & 207 & M2\\
IXa & 0.96 & 0.965 & 0.04 & 0.053 & 5--193 & 189 & M2\\
IXb & 0.965 & 0.975 & 0.045 & 0.053 & 5--150 & 146 & M2\\
\bottomrule\end{tabular}
\end{table}
\begin{table}[ht]
\centering\small
\caption{Trace-row ratio parameters.  In each row $D_f$ is given by
\eqref{bc:Dfloor}, so that the same lower bound is used in
\eqref{bc:kappa-gate}.}
\label{bc:trace-ratios}
\begin{tabular}{@{}crrrrrr@{}}\toprule
Row&Ia&Ib&II&III&IV&V\\\midrule
$\kappa$&2.04150406 & 2.11136288 & 2.0514852 & 2.04222636 & 2.04008696 & 2.0396935\\
\bottomrule\end{tabular}
\end{table}
\begin{table}[ht]
\centering\small
\caption{Targets and rational lower bounds for the final compact margins.}
\label{bc:compact-margins}
\begin{tabular}{@{}crr@{}}\toprule
Row&$T_0$&$\mathcal A(n_0)+\mathrm{bonus}-T_0>$\\\midrule
Ia & 25.574 & 193.4195\\
Ib & 11.111 & 8.6177\\
II & 24.88 & 44.1758\\
III & 24.378 & 16.8682\\
IV & 24.378 & 3.1044\\
V & 10.23 & 7.9635\\
VI & 10.116 & 5.4864\\
VII & 10.051 & .5545\\
VIIIa1 & 10.008 & .4089\\
VIIIa2 & 9.562 & .7142\\
VIIIb & 8.89 & 1.4735\\
IXa & 8.602 & 1.6317\\
IXb & 8.356 & .4620\\
\bottomrule\end{tabular}
\end{table}

The next two tables specify every contour choice by intervals of
integers.  The notation $[u,v]:k$ assigns the value $k$ to every
$j=u,u+1,\ldots,v$; a single integer denotes a singleton interval.
Assignments in a row are disjoint and cover exactly its $j$ interval
in Table~\ref{bc:compact-cover}.

\begin{longtable}{@{}lp{.84\textwidth}@{}}
\caption{The exact values of $10^4\nu_j$.}\label{bc:nu-table}\\
\toprule Row&Assignments\\\midrule\endfirsthead
\toprule Row&Assignments\\\midrule\endhead
Ia & $5:9975,\quad 6:9982,\quad 7:9987,\quad 8:9989,\quad 9:9991$\\
 & $10:9993,\quad [11,12]:9994,\quad [13,14]:9995,\quad [15,17]:9996,\quad [18,24]:9997$\\
 & $[25,257]:9998,\quad [258,264]:9997,\quad [265,266]:9996$\\
Ib & $6:9891,\quad 7:9919,\quad 8:9937,\quad 9:9950,\quad 10:9959$\\
 & $11:9965,\quad 12:9971,\quad 13:9975,\quad 14:9978,\quad 15:9980$\\
 & $16:9982,\quad 17:9984,\quad 18:9986,\quad 19:9987,\quad 20:9988$\\
 & $21:9989,\quad [22,23]:9990,\quad [24,25]:9991,\quad [26,27]:9992,\quad [28,30]:9993$\\
 & $[31,34]:9994,\quad [35,40]:9995,\quad [41,48]:9996,\quad [49,52]:9990,\quad [53,57]:9991$\\
 & $[58,64]:9992,\quad [65,68]:9993,\quad [69,70]:9986,\quad [71,75]:9987,\quad [76,78]:9988$\\
 & $[79,81]:9980,\quad [82,84]:9981,\quad 85:9971,\quad [86,88]:9972,\quad [89,90]:9961$\\
 & $91:9932$\\
II & $5:9967,\quad 6:9976,\quad 7:9982,\quad 8:9986,\quad 9:9989$\\
 & $10:9990,\quad 11:9992,\quad 12:9993,\quad [13,14]:9994,\quad [15,16]:9995$\\
 & $[17,20]:9996,\quad [21,29]:9997,\quad [30,188]:9998,\quad [189,206]:9997,\quad [207,211]:9996$\\
 & $212:9994$\\
III & $5:9980,\quad 6:9986,\quad 7:9989,\quad 8:9991,\quad 9:9993$\\
 & $10:9994,\quad [11,12]:9995,\quad [13,15]:9996,\quad [16,21]:9997,\quad [22,253]:9998$\\
 & $[254,262]:9997$\\
IV & $5:9984,\quad 6:9989,\quad 7:9991,\quad 8:9993,\quad 9:9994$\\
 & $10:9995,\quad [11,13]:9996,\quad [14,19]:9997,\quad [20,276]:9998,\quad [277,278]:9997$\\
V & $5:9985,\quad 6:9989,\quad 7:9992,\quad 8:9993,\quad 9:9994$\\
 & $10:9995,\quad [11,13]:9996,\quad [14,18]:9997,\quad [19,279]:9998,\quad [280,281]:9997$\\
VI & $5:9989,\quad 6:9992,\quad 7:9994,\quad 8:9995,\quad [9,10]:9996$\\
 & $[11,15]:9997,\quad [16,299]:9998$\\
VII & $5:9991,\quad 6:9993,\quad 7:9995,\quad [8,9]:9996,\quad [10,13]:9997$\\
 & $[14,244]:9998$\\
VIIIa1 & $5:9992,\quad 6:9994,\quad 7:9995,\quad [8,9]:9996,\quad [10,12]:9997$\\
 & $[13,248]:9998$\\
VIIIa2 & $5:9992,\quad 6:9994,\quad 7:9995,\quad [8,9]:9996,\quad [10,12]:9997$\\
 & $[13,236]:9998$\\
VIIIb & $5:9992,\quad 6:9994,\quad 7:9995,\quad [8,9]:9996,\quad [10,13]:9997$\\
 & $[14,211]:9998$\\
IXa & $5:9992,\quad 6:9994,\quad 7:9995,\quad 8:9996,\quad [9,12]:9997$\\
 & $[13,193]:9998$\\
IXb & $5:9992,\quad 6:9994,\quad 7:9995,\quad 8:9996,\quad [9,12]:9997$\\
 & $[13,150]:9998$\\
\bottomrule\end{longtable}

\begin{longtable}{@{}lp{.84\textwidth}@{}}
\caption{The exact values of $W_j$.}\label{bc:W-table}\\
\toprule Row&Assignments\\\midrule\endfirsthead
\toprule Row&Assignments\\\midrule\endhead
Ia & $[5,129]:1,\quad [130,184]:2,\quad [185,212]:3,\quad [213,233]:4,\quad [234,248]:5$\\
 & $[249,257]:6,\quad [258,264]:8,\quad [265,266]:10$\\
Ib & $[6,48]:1,\quad [49,68]:2,\quad [69,78]:3,\quad [79,84]:4,\quad [85,88]:5$\\
 & $[89,90]:6,\quad 91:8$\\
II & $[5,104]:1,\quad [105,149]:2,\quad [150,171]:3,\quad [172,188]:4,\quad [189,199]:5$\\
 & $[200,206]:6,\quad [207,211]:8,\quad 212:10$\\
III & $[5,127]:1,\quad [128,181]:2,\quad [182,209]:3,\quad [210,229]:4,\quad [230,244]:5$\\
 & $[245,253]:6,\quad [254,261]:8,\quad 262:10$\\
IV & $[5,134]:1,\quad [135,191]:2,\quad [192,220]:3,\quad [221,242]:4,\quad [243,257]:5$\\
 & $[258,268]:6,\quad [269,276]:8,\quad [277,278]:10$\\
V & $[5,135]:1,\quad [136,193]:2,\quad [194,222]:3,\quad [223,244]:4,\quad [245,260]:5$\\
 & $[261,270]:6,\quad [271,279]:8,\quad [280,281]:10$\\
VI & $[5,181]:1,\quad [182,253]:2,\quad [254,289]:3,\quad [290,299]:4$\\
VII & $[5,195]:1,\quad [196,244]:2$\\
VIIIa1 & $[5,217]:1,\quad [218,248]:2$\\
VIIIa2 & $[5,213]:1,\quad [214,236]:2$\\
VIIIb & $[5,204]:1,\quad [205,211]:2$\\
IXa & $[5,193]:1$\\
IXb & $[5,150]:1$\\
\bottomrule\end{longtable}

\subsection{Finite arithmetic verification}
\label{bc:arithmetic}

For the boundary calculations, all rational expressions in
Tables~\ref{bc:preliminary-table}--\ref{bc:endpoint-table} and
\eqref{bc:feedback} are evaluated with interval endpoints in $\mathbb Q$.
The following rules make the arithmetic finite and reproducible.
For $0\leq x\leq4$, put $S_{96}(x)=\sum_{j=0}^{96}x^j/j!$.
Then
\[
 S_{96}(x)<e^x<
 S_{96}(x)+\frac{x^{97}}{97!}\frac1{1-x/98}
 \quad(x>0);
\]
negative arguments are handled by reciprocal intervals.
For logarithms, first multiply the argument by a power of two to
place it in $[1,2)$.  If $t=(x-1)/(x+1)$ in that interval, then
\[
 2\sum_{j=0}^{95}\frac{t^{2j+1}}{2j+1}
 \leq\log x\leq
 2\sum_{j=0}^{95}\frac{t^{2j+1}}{2j+1}
 +\frac{2t^{193}}{193(1-t^2)}.
\]
Apply the same rule to $\log2$ when undoing the power-of-two change.
Square roots are enclosed between consecutive multiples of $10^{-56}$,
using integer square roots after multiplying the radicand by $10^{112}$.
For $\pi$, use
\[
 \pi=16\arctan(1/5)-4\arctan(1/239)
\]
with 100 alternating terms for each arctangent and the first omitted
term as its remainder bound.  These rules suffice to verify every
strict comparison in the boundary certificate.  Integer powers, even
those with exponent 1000, are evaluated exactly.

For the middle contours, only rational arithmetic, integer square
roots, and positive exponential and logarithmic sums are needed.
The choices for $e_W$ in \eqref{bc:inverse-gates-2}--\eqref{bc:inverse-gates-3}
are as follows:
\begin{center}\small
\begin{tabular}{@{}rrrr@{}}\toprule
$W$&$10^{12}e_W$&$W$&$10^{12}e_W$\\\midrule
1&367879441172&5&6737947000\\
2&135335283237&6&2478752177\\
3&49787068368&8&335462628\\
4&18315638889&10&45399930\\\bottomrule
\end{tabular}
\end{center}
The positive sums $S_{100}(W)=\sum_{j=0}^{100}W^j/j!$ satisfy
$S_{100}(W)>1/e_W$, and hence $e_W>e^{-W}$.
Likewise $S_{100}(T_0)>8/C_-(n_0)$ certifies the first comparison
in \eqref{bc:compact-final-gates}; all other terms in the final
budget are rational.

An explicit verification consists of the following finite steps.
For each of the thirteen rows, form $V,M,\alpha,\beta$ from
\eqref{bc:trace-field} or \eqref{bc:compact-moment-field} and verify
all its preceding field conditions.  For each integer $j$ in its
listed interval, read $\nu_j,W_j$ from the two assignment tables,
form \eqref{bc:inverse-def-1}--\eqref{bc:inverse-def-2}, and verify
\eqref{bc:inverse-gates-1}--\eqref{bc:inverse-gates-3},
$\nu_jH_0\widehat r_j>22/7$, and $k_j(n_0)>0$.
Next form \eqref{bc:Delta-def} and the outward radical
\eqref{bc:factor-bound}; verify $0<\Delta_j<1$, $0<\eta_j<1$,
and
\[
 (\eta_j-10^{-12})^2<
 \frac{B^2+R_j^2}{(1-\Delta_j)^2}\leq\eta_j^2.
\]
Compute the weight by \eqref{bc:weights}, and check its positivity,
its nonincrease with $j$, and the strict increase of $k_j(n_0)$.
Finally verify the two equivalent forms of the Abel sum in
\eqref{bc:Abel}, the positive slope and intercept loss $A_*,B_*$,
and the strict margins in Table~\ref{bc:compact-margins}.
There are exactly 2918 such contours, all specified by the printed
tables.

The two bridge contours use $(\nu,W,e_W)=(.99,6,1/400)$ with the
parameters given in Proposition~\ref{bc:bridge}.  Their remaining
checks are \eqref{bc:inverse-gates-1}--\eqref{bc:inverse-gates-3},
the displayed factor and count bounds, and the scalar comparisons
in its four cases.  Upper bounds for $e^x$ at
$x=1,.75,.74$ follow from
\[
 e^x<\sum_{j=0}^{24}\frac{x^j}{j!}
       +\frac{x^{25}}{25!}\frac1{1-x/26};
\]
all the other bridge exponential comparisons use $S_{100}$ from below.
The high-modulus radical tests are the two positive-side squared
comparisons following \eqref{bc:high-necessary}.
Thus the entire compact argument reduces to the displayed finite
integer and rational operations, followed by the explicit monotonicity
arguments for $n\geq6251$.
